\documentclass{article}
\usepackage[utf8]{inputenc}
\usepackage{xspace}
\usepackage[top=1.5cm, bottom=1.5cm, left=1.5cm, right=1.5cm]{geometry}
\usepackage{enumitem}
\usepackage{graphicx}
\usepackage{textcomp}
\usepackage{amssymb}
\usepackage{ifsym}
\usepackage{xspace}
\usepackage{amsmath}

\title{Facultad de Ciencias. \\Bitácora 1 \\ Geoquímica orgánica}
\author{Arteaga Hernández Alejandro Iabin }
\date{Entrega 7 de octubre 2020}

\begin{document}

\maketitle

\begin{enumerate}
    \item ¿Qué es un análisis cualitativo y análisis cuantitativo? ¿cuál es la
diferencia entre ellos?

En el Análisis Cualitativo Orgánico, se pone atención a la identificación de los elementos y grupos funcionales que integran la muestra. Dado que se trata de materia orgánica, las estructuras se vuelven en ocasiones muy complejas, y la sistematización más difícil.


El Análisis Cuantitativo está basado en la aplicación de las leyes de la Estequiometria. Se procede tomando una cantidad bien determinada de muestra, con peso o volumen conocido, y sometiéndola a reacciones químicas


 \item ¿Qué es una reacción REDOX? ¿cuáles son los estados de
oxidación del carbono e hidrógeno?

Redox es el nombre que recibe una reacción de tipo químico que implica la transferencia de electrones entre distintos reactivos, lo que lleva a una modificación del estado de oxidación. En estas reacciones, un elemento pierde electrones y el otro, los recibe.


Carbono: Estado de oxidación	+4

Hidrógeno: estado de oxidación +1 , -1

\item ¿Cuál es la diferencia entre un mechero de Bunsen y una mufla
si ambos son fuentes de calentamiento?


La mufla tiene un  volumen de trabajar mayor a demás de alcanzar temperaturas de hasta 1100°c y se utiliza principalmente para calcinar muestras , mientras que el mechero Bunsen se utiliza como fuente de energía calorífica por su tamaño portátil y las temperaturas están limitadas por la llama del gas utilizado


La mufla debe ser fija, y recubierta y usa electricidad, el mechero usa gas.

\item ¿Cuáles son los diagramas de la NFPA de los siguientes
compuestos: 

$CaCO_3, BaCO_3, Ba(OH)_2, CuSO_4, CuSO_{4(s)} · 5H2O
y CO_2$?
\end{enumerate}

\section{Fuentes:}

https://www.ejemplode.com/38-quimica/4746-analisis\_cualitativo\_y\_cuantitativo.html
\\\\
https://definicion.de/redox/

\end{document}
